\documentclass[11pt]{letter}
\usepackage[margin=1in]{geometry}
\signature{A.~Author\\[Institution]}
\address{[Institution]\\[City]}
\date{\today}
\begin{document}
\begin{letter}{Editors\\Physical Review Letters}
\opening{Dear Editors,}

Please consider the manuscript ``Nucleation density programs the viscoelastic gel point of a percolation network'' for publication as a Letter. There is no related submission history at Physical Review, and the manuscript is not part of a joint submission. We have no recommended or excluded referees.

\textbf{Justification.}
Passive microrheology already locates the gel point, and growth-rule percolation already moves a threshold. This Letter joins those precedents. A probe in the void reads the gel point as the void-spanning transition, and one nucleation density shifts it from the Bernoulli value toward unity. Finite-size scaling gives one effective exponent, about 1.6, with no trend in that density and distinct from the static value 0.88. Rheologists gain a morphological handle on where a gel sets, and the percolation community gains a dynamically defined threshold. The result is sharp enough, and the audience broad enough, for a Letter.

The finite-size protocol and the three-dimensional walk-dimension argument ($2/d_w\approx0.52$) are in the End Matter. The GSER spectra, the full nucleation-density table, and the code notes are in the Supplemental Material. Reference titles are included. Data and code will be deposited at Zenodo; the data-availability statement in the manuscript names that deposit and will carry the DOI at publication.

\closing{Sincerely,}
\end{letter}
\end{document}
